\chapter[CONCLUSION AND FUTURE WORK]{Conclusion and Future Work}\label{ch:V}

\section{Conclusion}\label{sec:5.1}

Taken together, the two testing groups suggested that the application performed
strongly on its core function, indoor navigation, while leaving room for
refinement in the layout builder. Users rated the navigation experience highly,
with a mean score falling within the ``Excellent'' range. Building
administrators, on the other hand, found the layout builder and sandbox mode
useful for building layout planning, helping communicate layout proposals to
stakeholders and clients. Some limitations of the application evaluation include
the small sample size of both testing groups, consisting of six end users and
four administrators. This limits how confidently these results can be
generalized to a wider range of users. Additionally, testing was limited to a
one-week period for both types of participants, which does not capture how
usability holds up with repeated or long-term use of the application. This is
particularly relevant for building administrators, who may create and manage
large, multi-story buildings over time.

\section{Future Work}\label{sec:5.2}

For future work, a primary area of focus is expanding the system beyond its
current scope to cover other buildings across the campus, allowing users to
navigate between buildings rather than within a single building alone. To
complement this expansion, the development of native iOS and Android
applications is recommended, as a transition from a web-based platform to
dedicated mobile applications could offer an improved overall user experience.

The application should also integrate accessibility features, which involve
mapping and displaying accessible routes, such as paths incorporating PWD
(persons with disability) ramps, elevators, and accessible restrooms. Relatedly,
integrating clear displays of emergency exit routes is recommended to enhance
the safety and preparedness of all users within the building.

With respect to the navigation feature specifically, future versions should
support real-time location tracking to guide users continuously as they move
toward their destination, rather than relying on a static route display, and
should display doorways along the suggested path to further improve route
clarity. Further routing optimization and testing are also recommended to
broaden the scope of navigation scenarios evaluated, including more varied
stairway configurations and larger, multi-floor paths.

On the administrator side, the layout builder would benefit from several
refinements identified during testing. These include an autosave feature to
prevent the loss of in-progress map drafts, a clearer distinction between the
Layout Builder and Sandbox features to reduce confusion between the two, and
support for real-world unit scaling with gridlines or rulers to improve drawing
precision. The ability to upload a 2D or CAD reference plan to serve as a guide
when constructing more accurate layouts is also recommended. Improving the
responsiveness and synchronization of the layout builder's editing tools would
further support administrators working on larger, more complex building layouts.

To support long-term map accuracy, a crowdsourced update and feedback system is
recommended, through which users can report inaccuracies or suggest corrections
directly within the application, helping keep building layouts current without
relying solely on administrator updates.

Finally, integrating the system with outdoor campus navigation is recommended to
provide a connected wayfinding experience for users moving between buildings.
